\subsubsection{Rozložení a navigace}

Požadavky na rozložení a navigaci jsou zamřeney na to, jak by se uživatel měl v systému orientovat \parencite{rejnkovaLokalizacePrizpusobeniMetodiky}. Tyto požadavky jsou popsány v tabulce \ref{tab:ROZ-pozadavky}.

\begin{OstatniPozadavkyTabulka}{ROZ}{Požadavky na rozložení a navigaci}
  \OstatniPozadavekRow{Mapa jako hlavní rozcestník}{Hlavní stránka veřejného uživatelského rozhraní musí zobrazovat interaktivní mapu se všemi registrovanými budovami. Mapa slouží jako hlavní rozcestník pro vyhledávání a navigaci k detailu jednotlivých entit.}{1}
  \OstatniPozadavekRow{Drobečková navigace v hierarchii entit}{V detailu místnosti, senzoru a měření musí být dostupná drobečková navigace (\emph{breadcrumb}) zobrazující cestu hierarchií budova $\rightarrow$ místnost $\rightarrow$ senzor $\rightarrow$ měření. Každý prvek navigace je klikatelný a vede na detail nadřazené entity.}{2}
  \OstatniPozadavekRow{Časový graf měření na detailu senzoru}{Detail senzoru musí obsahovat časový graf zobrazující průběh měření v čase pro každou veličinu, kterou senzor měří. Uživatel musí mít možnost zvolit časový rozsah (posledních 24 h, 7 dní, 30 dní, vlastní rozsah).}{1}
\end{OstatniPozadavkyTabulka}
